\section{Related Work}
\label{sec:related-work}

\reflector{} is best understood as a synthesis across several established traditions rather than as a standalone invention. The relevant background includes cybernetic feedback and regulation, bounded rationality, mixed-initiative and human-in-the-loop AI, distributed cognition, distributed synchronization theory, and software workflow orchestration. This section therefore takes a selective positioning approach: it identifies the conceptual and technical lineages most relevant to \reflector{}'s architecture, then clarifies where \reflector{} combines them in a distinct way.

\subsection{Cybernetics and Feedback Regulation}

Cybernetics frames intelligent behavior as control under feedback: systems regulate action by comparing observed state to target state and adapting accordingly \cite{Wiener1948Cybernetics}. Later systems work emphasized multi-level organizational control, where viability depends on recursively coupled regulation layers rather than on a single command point \cite{Beer1972BrainFirm}. \reflector{} extends this feedback logic into AI-assisted engineering workflows. Reflective audits, synchronization checkpoints, and continuation gates play a cybernetic role: they periodically re-couple execution with intent so recursive activity remains governable over time.

\subsection{Bounded Rationality and Cognitive Constraints}

Bounded rationality research shows that decision-makers operate under limited attention, incomplete information, and finite time \cite{Simon1957ModelsMan}. Cognitive load theory similarly highlights that coordination and interpretation overhead can become a primary bottleneck in complex problem solving \cite{Sweller1988CognitiveLoad}. \reflector{} adopts these constraints as design premises. Its scoped delegation, checkpointing, and artifact reconciliation mechanisms are intended to reduce coordination overload by compressing recursive state into inspectable synchronization moments, rather than requiring continuous human reconstruction of every intermediate step.

\subsection{Mixed-Initiative and Human-in-the-Loop Systems}

Interactive machine learning and human-centered AI literature argues that reliable collaborative systems emerge when humans and automation share initiative, with humans retaining authority over normative and contextual judgments \cite{Amershi2014InteractiveML,Shneiderman2020HCAI}. Recent software-engineering agent systems reinforce this framing: language-model-driven execution can be high-throughput, but practical reliability remains tightly coupled to task framing, artifact quality, and evaluation context \cite{Chen2021Codex,Jimenez2024SWEBench,Yang2024SWEAgent}. \reflector{} aligns with this tradition by treating human oversight as structural, not exceptional: recursive continuation is explicitly conditioned on synchronized evidence and human-reviewed boundaries.

\subsection{Distributed Cognition and Artifact-Mediated Coordination}

Distributed cognition treats cognition as spread across people, tools, representations, and environments rather than confined to individual minds \cite{Hutchins1995CognitionWild}. \reflector{} adopts this view directly: issues, specifications, pull requests, audit traces, and repository history are treated as shared cognitive infrastructure for recursive work. In this framing, synchronization is not only process control; it is also cognitive state maintenance across participants and artifacts that evolve asynchronously.

\subsection{Synchronization, Observability, and Global State}

Distributed systems research formalizes synchronization as the challenge of establishing meaningful ordering and recoverable global state under partial views \cite{Lamport1978TimeClocks,Chandy1985DistributedSnapshots}. Work on eventual consistency highlights practical tradeoffs between local progress and globally coherent state \cite{Vogels2009EventuallyConsistent}. \reflector{} borrows these intuitions at the workflow level. Synchronization checkpoints act as bounded global-state reconstructions for recursive engineering activity, while reflective audits provide observability surfaces needed to decide whether continuation is warranted, deferred, or redirected.

\subsection{Developer Workflow Orchestration and Platform Framing}

Software engineering and architecture literature has long emphasized that structure, modular decomposition, and coordination strategy shape system outcomes as much as raw implementation effort \cite{Parnas1972Modularization,Brooks1987NoSilverBullet}. \reflector{} intersects with contemporary developer-experience and platform concerns in this same spirit: it treats workflow design, artifact interfaces, and governance boundaries as first-class infrastructure. Unlike pure CI/CD pipelines or productivity tooling, however, its primary objective is recursive synchronization stability under mixed-initiative execution.

\subsection{\reflector{}'s Distinct Synthesis}

The central claim is intentionally narrow. \reflector{} is not presented as a general autonomous agent framework, a conventional coding assistant, or a replacement for existing software lifecycle systems. Its contribution is the integration of: (i) cybernetic feedback principles, (ii) bounded-rational coordination assumptions, (iii) mixed-initiative governance, and (iv) synchronization/observability mechanisms into a single recursive orchestration architecture for AI-assisted software engineering.

In other words, prior work supplies many of \reflector{}'s components, but usually in isolation: feedback and regulation traditions explain stabilization, mixed-initiative work explains collaborative authority, distributed cognition explains artifact-mediated reasoning, and distributed synchronization theory explains reconciliation under partial state. \reflector{}'s proposed novelty is the operational combination of these components as recursive stabilization infrastructure for real engineering workflows.

This positioning is meant to be intellectually conservative: \reflector{} depends on established ideas while recombining them around a specific systems problem---maintaining coherence, interpretability, and governance as recursive AI-assisted execution scales. The next section extends this positioning into a forward-looking research agenda (\Cref{sec:future-directions}).
